% ============================================================
%  Part XIII — Reference lists
% ============================================================
\gpart{XIII}{Reference}

% ------------------------------------------------------------
\gsec{Verbs that take the dative}

\gintro{A group of common verbs takes a dative object where English takes a direct one. There is no pattern to them, so they are listed here in full.}

These verbs have no accusative object at all. Their single object is dative,
even though English treats it as a direct object.

\begin{gtbl}{colspec={X[0.8,l] X[1,l] X[0.8,l] X[1,l]}}
  Verb & Meaning & Verb & Meaning \\
  helfen     & to help        & danken     & to thank \\
  folgen     & to follow      & gehören    & to belong to \\
  gefallen   & to please      & passen     & to suit, fit \\
  antworten  & to answer      & glauben    & to believe \eng{a person} \\
  gratulieren & to congratulate & begegnen & to meet by chance \\
  schmecken  & to taste good to & fehlen   & to be missing to \\
  raten      & to advise      & vertrauen  & to trust \\
  zuhören    & to listen to   & ähneln     & to resemble \\
  dienen     & to serve       & drohen     & to threaten \\
  imponieren & to impress     & nützen     & to be of use to \\
  schaden    & to harm        & wehtun     & to hurt \\
  widersprechen & to contradict & verzeihen & to forgive \\
\end{gtbl}

\begin{gnote}
\textit{Ich helfe \textbf{dir}}, not \textit{ich helfe dich}. A useful test:
these verbs cannot be turned into a normal passive, precisely because they have
no accusative object to promote.
\end{gnote}

\tcap{Dative and accusative together}

\begin{flattbl}{colspec={X[0.75,l] X[2.25,l]}}
  geben     & Ich gebe \textbf{dir} \textbf{das Buch}. \\
  zeigen    & Er zeigt \textbf{mir} \textbf{die Stadt}. \\
  schenken  & Sie schenkt \textbf{ihm} \textbf{eine Uhr}. \\
  erklären  & Ich erkläre \textbf{euch} \textbf{die Regel}. \\
  schicken  & Wir schicken \textbf{ihnen} \textbf{einen Brief}. \\
  Others    & bringen, kaufen, sagen, empfehlen, anbieten, leihen \\
\end{flattbl}

% ------------------------------------------------------------
\gsec{Verbs with fixed prepositions}

\gintro{Many verbs are locked to a particular preposition, and that preposition then fixes the case. English does the same thing (\emph{to wait \textbf{for}}), but the pairings rarely match.}

The preposition is part of the verb and its case is fixed, so the verb, the
preposition and the case work as a single unit.

\begin{gtbl}{colspec={X[1.15,l] X[0.6,c] X[1.25,l]}}
  Verb + preposition & Case & Meaning \\
  warten auf         & Akk. & to wait for \\
  sich freuen auf    & Akk. & to look forward to \\
  sich freuen über   & Akk. & to be pleased about \\
  denken an          & Akk. & to think of \\
  sich erinnern an   & Akk. & to remember \\
  sich interessieren für & Akk. & to be interested in \\
  sich kümmern um    & Akk. & to take care of \\
  sprechen über      & Akk. & to talk about \\
  sich ärgern über   & Akk. & to be annoyed about \\
  bitten um          & Akk. & to ask for \\
  glauben an         & Akk. & to believe in \\
  sich gewöhnen an   & Akk. & to get used to \\
  fragen nach        & Dat. & to ask about \\
  suchen nach        & Dat. & to search for \\
  sich treffen mit   & Dat. & to meet with \\
  teilnehmen an      & Dat. & to take part in \\
  leiden unter       & Dat. & to suffer from \\
  sich fürchten vor  & Dat. & to be afraid of \\
  abhängen von       & Dat. & to depend on \\
  sprechen mit       & Dat. & to speak with \\
\end{gtbl}

\begin{gnote}
Answer these with a \textit{da(r)-} compound rather than a pronoun when the
object is a thing: \textit{Ich warte \textbf{darauf}}, \textit{Ich denke
\textbf{daran}}. Keep the preposition plus pronoun for people:
\textit{Ich warte \textbf{auf ihn}}.
\end{gnote}

% ------------------------------------------------------------
\gsec{Numbers}

\gintro{German numbers are written as one word and say the units before the tens, which is the main thing to get used to.}

\begin{gtbl}{colspec={X[0.5,c] X[1,l] X[0.5,c] X[1,l] X[0.5,c] X[1,l]}}
  \# & Wort & \# & Wort & \# & Wort \\
  0 & null   & 10 & zehn     & 20 & zwanzig \\
  1 & eins   & 11 & elf      & 21 & einundzwanzig \\
  2 & zwei   & 12 & zwölf    & 30 & dreißig \\
  3 & drei   & 13 & dreizehn & 40 & vierzig \\
  4 & vier   & 14 & vierzehn & 50 & fünfzig \\
  5 & fünf   & 15 & fünfzehn & 60 & sechzig \\
  6 & sechs  & 16 & sechzehn & 70 & siebzig \\
  7 & sieben & 17 & siebzehn & 80 & achtzig \\
  8 & acht   & 18 & achtzehn & 90 & neunzig \\
  9 & neun   & 19 & neunzehn & 100 & (ein)hundert \\
\end{gtbl}

\tcap{Building bigger numbers}

\begin{flattbl}{colspec={X[0.8,l] X[2.2,l]}}
  Units first & 21 = \textit{ein\textbf{und}zwanzig}, 47 = \textit{sieben\textbf{und}vierzig} \\
  One word    & 365 = \textit{dreihundertfünfundsechzig} \\
  Irregular   & \textit{dreißig} \eng{-ßig}, \textit{sechzig}, \textit{siebzig} \eng{shortened} \\
  Big         & 1\,000 = \textit{tausend}, 1\,000\,000 = \textit{eine Million} \\
  Years       & 1999 = \textit{neunzehnhundertneunundneunzig}; 2026 = \textit{zweitausendsechsundzwanzig} \\
\end{flattbl}

\begin{gnote}
German swaps the decimal point and the thousands separator:
\textbf{1.000,50} is one thousand point five zero. \textit{eins} loses its
\textit{-s} before a noun: \textit{ein Buch}, but \textit{Nummer eins}.
\end{gnote}

% ------------------------------------------------------------
\gsec{Time and date}

\gintro{Clock times, days and dates follow a few fixed patterns. Most of them involve a preposition contracted with the article.}

\tcap{Telling the time}

\begin{gtbl}{colspec={X[0.7,l] X[1.15,l] X[1.15,l]}}
  Clock & Umgangssprache & Formal \\
  8:00  & acht Uhr             & acht Uhr \\
  8:15  & Viertel nach acht    & acht Uhr fünfzehn \\
  8:30  & \textbf{halb neun}   & acht Uhr dreißig \\
  8:45  & Viertel vor neun     & acht Uhr fünfundvierzig \\
  8:20  & zwanzig nach acht    & acht Uhr zwanzig \\
  8:40  & zwanzig vor neun     & acht Uhr vierzig \\
\end{gtbl}

\begin{gnote}
\textbf{halb neun} means \emph{half way to nine}, i.e. 8:30 — not half past
nine. It is a common source of confusion for English speakers.
\end{gnote}

\tcap{Days and months}

\begin{flattbl}{colspec={X[1,l] X[1,l] X[1,l] X[1,l]}}
  Montag  & Dienstag & Mittwoch & Donnerstag \\
  Freitag & Samstag  & Sonntag  & \eng{all der} \\
  Januar  & Februar  & März     & April \\
  Mai     & Juni     & Juli     & August \\
  September & Oktober & November & Dezember \\
\end{flattbl}

\tcap{Dates}

\begin{flattbl}{colspec={X[0.8,l] X[2.2,l]}}
  Writing  & \textit{der 1.\ Mai} \eng{read: der erste Mai} \\
  On a day & \textbf{am} Montag, \textbf{am} 1.\ Mai \\
  In a month & \textbf{im} Mai, \textbf{im} Sommer \\
  In a year & \textit{2026} or \textit{im Jahr 2026} — never \textit{in 2026} \\
  Letters  & \textit{Berlin, den 1.\ Mai 2026} \\
\end{flattbl}

% ------------------------------------------------------------
